\section{Supplementary Methods and Results}
\label{sec:supplementary}

\subsection{Audit-preserving annotation workflow}

The RQ6 review package assigned a stable anonymous passage identifier within each separately randomized workbook. A private mapping connected that identifier to the frozen source chunk. Reviewer files did not expose method provenance or Gold anchors. The audit archive preserves the original Reviewer A and Reviewer B workbooks, a normalized machine-readable label ledger, the disagreement list, adjudication notes, and the final 58-query run package. Adjudication produced a new decision layer and did not overwrite either reviewer's submission.

\begin{table}[htbp]
\centering
\begin{threeparttable}
\caption{Agreement before joint adjudication.}
\label{tab:agreement}
\begin{tabular}{lccc}
\toprule
Label level & Exact agreement & Cohen's $\kappa$ & Gwet AC1 \\
\midrule
Question status & 0.967 & 0.000 & 0.966 \\
Evidence completeness & 0.967 & 0.000 & 0.966 \\
Five-class passage label & 0.663 & 0.488 & 0.599 \\
Core binary Gold & 0.941 & 0.792 & 0.918 \\
\bottomrule
\end{tabular}
\begin{tablenotes}[flushleft]\footnotesize
\item Question-level $\kappa=0$ reflects extreme marginal imbalance: Reviewer A used one category for all questions, while Reviewer B identified two incomplete questions. Core Gold merges \texttt{Sufficient} and \texttt{Required Component}.
\end{tablenotes}
\end{threeparttable}
\end{table}


The question-level $\kappa$ values in \Cref{tab:agreement} should not be read as zero practical agreement. Both reviewers agreed on 58 of 60 questions, but one reviewer's use of a single category makes the expected-agreement term degenerate. The passage-level core Gold comparison has non-degenerate positive and negative classes and is the agreement result relevant to retrieval scoring.

\subsection{Selective reranking confirmation}

\begin{table}[htbp]
\centering
\begin{threeparttable}
\caption{Selective reranking confirmation set ($n=30$; query-bootstrap 95\% confidence intervals).}
\label{tab:rq5}
\small
\begin{tabularx}{\textwidth}{l*{3}{>{\centering\arraybackslash}X}}
\toprule
Method & Hit@5 & MRR & nDCG@5 \\
\midrule
Context-matched BM25 & \textbf{0.900 [0.800, 1.000]} & 0.798 [0.667, 0.917] & 0.785 [0.656, 0.901] \\
R1 source-chunk dense & 0.800 [0.633, 0.933] & 0.718 [0.575, 0.853] & 0.716 [0.571, 0.850] \\
Unconditional source RRF & 0.833 [0.700, 0.967] & \textbf{0.832 [0.699, 0.941]} & \textbf{0.798 [0.655, 0.919]} \\
Selective table gate & 0.867 [0.733, 0.967] & 0.790 [0.653, 0.912] & 0.774 [0.636, 0.897] \\
\bottomrule
\end{tabularx}
\begin{tablenotes}[flushleft]\footnotesize
\item The selective gate did not satisfy the prespecified MRR superiority criterion.
\end{tablenotes}
\end{threeparttable}
\end{table}


The selective gate activated on ten RQ5 queries: five through explicit table language and five through candidate-table support. It improved reciprocal rank on one query, reduced it on one, and left eight unchanged. In the failure case, BM25 ranked a Gold chunk first, while that chunk did not appear in the dense top 30; fusion moved it to rank 13. This example motivated the later decision to rerank a fixed BM25 candidate set rather than weakening lexical candidate recall.

\subsection{Locked Qwen3 configuration}

The supplementary neural experiment used \texttt{Qwen/Qwen3-Reranker-0.6B} from the official ModelScope repository. Nine model files, including the 1,191,588,280-byte \texttt{model.safetensors}, matched the repository SHA-256 values before use. The complete model-file manifest has SHA-256 \texttt{9e345ac4f295b7fe6f9734e1d5c1c73625b7a01ef2217be1896bb608c13fa508}.

The lock specifies BM25 candidate depth 50, maximum model length 1,024, BF16 inference, batch size eight, no fine-tuning, deterministic tie-breaking by BM25 rank and chunk identifier, and bootstrap seed 20260716. The lock validation passed candidate, corpus, Gold-package, model-manifest, and code-hash checks before inference. The complete inference retained 50 unique candidates for all 58 queries. The result-validation artifact has SHA-256 \texttt{4fabee1e7af0fca525ee23f307ffaf634052a6928419ec883b523c43efcbc565}.

\begin{table}[htbp]
\centering
\begin{threeparttable}
\caption{Supplementary Qwen3 and BM25 results by query slice.}
\label{tab:qwen_slices}
\begin{tabular}{llccc}
\toprule
Slice ($n$) & Method & Hit@5 & MRR@50 & nDCG@5 \\
\midrule
Single clause (20) & BM25 & 0.950 & 0.806 & 0.840 \\
 & Qwen3 & 0.950 & 0.844 & 0.846 \\
Document structure (15) & BM25 & 1.000 & 0.802 & 0.803 \\
 & Qwen3 & 1.000 & 0.889 & 0.864 \\
Table (13) & BM25 & 0.923 & 0.846 & 0.856 \\
 & Qwen3 & 1.000 & 0.827 & 0.871 \\
Cross-document (10) & BM25 & 0.900 & 0.842 & 0.644 \\
 & Qwen3 & 1.000 & 0.728 & 0.619 \\
\bottomrule
\end{tabular}
\begin{tablenotes}[flushleft]\footnotesize
\item Slice estimates are descriptive. The Qwen3 analysis is supplementary and post hoc.
\end{tablenotes}
\end{threeparttable}
\end{table}


The reranker changed the distribution of early ranks without changing candidate recall. It improved document-structure ranking most clearly in this small pool, while reducing cross-document reciprocal rank. These mixed effects are why the manuscript reports the aggregate positive signal without a general superiority claim.

\subsection{Reproducibility boundaries}

The public repository contains the experiment code and a text-free lock manifest. Candidate files and review materials contain source passages and remain outside the repository pending licensing review. Local audit paths are machine-specific and are not presented as public data access. The package-level citation audit records the primary or authoritative page used to verify every bibliography entry and the manuscript claim it supports.
